\begin{theorem}[The gather-based composition]
  \label{thm:aba-afw-protocol-composed}
  \lean{PLTS.ABA.AFW.broadcastReturnsFor}\lean{PLTS.ABA.AFW.roundProjection}\lean{PLTS.ABA.AFW.BroadcastReturnsInvariant}\lean{PLTS.ABA.AFW.ProtocolRelation}\lean{PLTS.ABA.AFW.protocolSimulation}\lean{PLTS.ABA.AFW.protocol_composed}\leanok
  \uses{def:aba-afw-protocol, def:aba-afw-round-families, thm:forward-sound, def:achievable}
  The gather-based protocol is forward-simulated by its composed system, and hence every trace distribution it achieves is achieved there.
  The relation has six conjuncts: the round loops and the coin oracle are shared objects, the ABA network is the DECIDED sets beside the corrupted set, each round's instance is the view $\mathrm{roundProjection}$ of the implementation's state, the bound bit of every round whose second gather has been called is on record, and the broadcast invariant holds at every broadcast instance of the view.
  The first four are unguarded, so they determine the composed state from the implementation's; the fifth, \leandecl{PLTS.ABA.AFW.BoundInvariant}, and the sixth, \leandecl{PLTS.ABA.AFW.BroadcastReturnsInvariant}, are the two clauses that are not read off the implementation's state.
  Assembling the view undoes the three separations between the implementation and the composed system --- the local states are transposed back out of the round records the $n$ processes hold, each network's sent set is projected out of the adversary's single tagged sent-set family, and what each instance returned to a program is read off its own local states in the broadcast instances as the values with a vote quorum there, \leandecl{PLTS.ABA.AFW.broadcastReturnsFor} --- and it reads the two gathers' cores and the round's bound bit off the adversary's ghost record of that round.
  The broadcast invariant is what identifies a returned value with the value a guard of the implementation names: a row of the implementation hands its composed counterpart a specific value with a vote quorum, the view chose that value, and under the invariant two vote quorums at one process name one value.
\end{theorem}

\begin{proof}\leanok
  \uses{def:aba-afw-protocol, def:aba-afw-round-families, def:forwardsim}
  Every row writes the acting process's round record and records at most one tagged message, and one lemma says what that does to the view: the round the row names reads as the one-point update of each of its coordinates.
  What a row then owes is sent set algebra --- which of the six slices takes the message and which are left alone --- and slicing commutes with a sent set write in exactly those two ways.
  On that basis each row of the implementation is answered by a run of the composed round: a send and a delivery are hidden events of the round, answered by one silent step, the adversary's authenticity conjunct becoming membership in the projected sent set; the graded-agreement call is the round's own call; and the three fused rows (D28) are answered by two steps through a named intermediate state --- the return-then-call row by the first gather's return and the second gather's call, the graded return by the second gather's return and the round's own, and a delivery completing a vote quorum by the instance's delivery and its return, which records what the instance returned.
  The invariant is re-established after every matched run by the broadcast instance's own preservation lemma, every matched step being a genuine step of the instance.
  The remaining labels never reach the implementation.
  The ABA interface, the coin handshake, the DECIDED rows and the fused coin return move the round loop and the ABA network while the family of rounds stands still; corruption is one broadcast, the corrupted set the adversary holds being the corrupted set of every network state under the same guard; and the two Byzantine graded-agreement rows have no row at the process they name.
  A program's row on a label of $\mathrm{roundOwn}(j)$ is a row of the implementation, which is what lets each of those cases rule the others out.
  The three fields of the round instance that no guard of it reads --- the core of each of its two gather instances and the round's bound bit --- are the adversary's ghost record at the implementation, and a row's ghost write is the instance's write of them.
  The view of the round the row names is the one-point update of those three coordinates, \leandecl{PLTS.ABA.AFW.roundProjection_writeGhost}; every other round's view stands where it is, \leandecl{PLTS.ABA.AFW.roundProjection_writeGhost_ne}; and a row whose write returns the record it found moves no round's view at all, \leandecl{PLTS.ABA.AFW.roundProjection_ghostId}.
  The return-then-call row writes the first gather's core and the bound bit, and it is the only row that extends the fifth conjunct, so the bit a graded return announces is on record when the return fires; the composed round's return drops the recorded grade, so the view's program records hold no grade outside a run.
\end{proof}
